\documentclass[11pt]{article}
\usepackage[a4paper,margin=27mm]{geometry}
\usepackage{amsmath,amssymb,amsthm,xcolor,booktabs,array}
\newcommand{\PROVED}{\textcolor{green!40!black}{\textbf{PROVED}}}
\newcommand{\OPEN}{\textcolor{red!70!black}{\textbf{OPEN}}}
\newcommand{\AUDIT}{\textcolor{orange!70!black}{\textbf{AUDIT}}}
\newtheorem{theorem}{Theorem}
\newtheorem{proposition}{Proposition}
\title{Poisson--Tate audit of the half-Tate Banach descent (v102)}
\date{13 August 2026}

\begin{document}
\maketitle

\section{Mellin form of the two orientations}

For $a=(a_n)\in\mathcal B_{1/2}$ put
\[
 A_af(x)=\sum_{n\ge1}a_nf(nx),\qquad
 D_a(s)=\sum_{n\ge1}a_nn^{-s}.
\]
With $\mathscr Jf(x)=x^{-1}f(x^{-1})$, direct calculation gives
\begin{align}
 \widehat{A_af}(s)&=D_a(s)\widehat f(s),             \tag{1}\\
 \widehat{\mathscr JA_a\mathscr Jf}(s)
     &=D_a(1-s)\widehat f(s).                        \tag{2}
\end{align}
Thus the Poisson-symmetric sum has multiplier
\begin{equation}
                         D_a(s)+D_a(1-s).             \tag{3}
\end{equation}

\begin{proof}
For (1), substitute $y=nx$ in the Mellin integral.  Moreover
$\mathscr JU_n\mathscr J=n^{-1}U_{1/n}$; its Mellin eigenvalue is
$n^{s-1}$.  Summing gives (2).
\end{proof}

\section{A two-pole test family}

Fix $0<\varepsilon<1/2$ and take
\begin{equation}
                         a_n=n^{-1/2-\varepsilon}.    \tag{4}
\end{equation}
Then
\[
 \|a\|_{1/2}=\sum_{n\ge1}n^{-1-\varepsilon}<\infty,
 \qquad D_a(s)=\zeta(s+1/2+\varepsilon).             \tag{5}
\]
The two terms in (3) have poles at
\begin{equation}
 s_-=\frac12-\varepsilon,
 \qquad s_+=\frac12+\varepsilon.                    \tag{6}
\end{equation}

\begin{theorem}[Two Tate modes do not renormalize the Banach action ---
\PROVED]
Choose $f\in\mathcal H_-$ such that
$\widehat f(s_-)\widehat f(s_+)\ne0$.  Mellin inversion of the two
oriented action produces nonzero power asymptotics proportional to
\begin{equation}
                 x^{-s_-}\quad\text{and}\quad x^{-s_+}.        \tag{7}
\end{equation}
Neither exponent is a Tate exponent $0$ or $1$.  Consequently subtracting
the two polar/Tate modes cannot make the symmetrized half-Tate Banach action
preserve $\mathcal H_-$.  This remains true before any question about the
zeros of $\zeta$.
\end{theorem}

\begin{proof}
Equations (3)--(6) show that the Mellin transform of the symmetrized output
has simple poles at $s_-$ and $s_+$ with respective nonzero residues
$\widehat f(s_-)$ and $\widehat f(s_+)$.  Moving the Mellin inversion line
across these poles yields their residue terms (7).  Elements of
$\mathcal H_-=\mathcal S_{(-\infty,\infty)}$ decay faster than every power
at both ends, so either nonzero term excludes membership in $\mathcal H_-$.
The even quotient of row (c) removes only the characters $s=0$ and $s=1$;
because $0<s_-<s_+<1$, those two corrections cannot cancel (7).
\end{proof}

The test function exists because the two Mellin evaluations are distinct
continuous nonzero functionals on $\mathcal H_-$; the union of their two
kernels is not the whole vector space.

\section{Consequence for the proposed descent}

Theorem 1 rules out the following route:
\[
 \mathcal B_{1/2}\curvearrowright\mathcal S_>
 \quad\xrightarrow{\text{Poisson orientation + two Tate subtractions}}\quad
 \mathcal B_{1/2}\curvearrowright\mathcal H_-^0.
\]
The failure is not an uncontrolled remainder: it is the explicit pair
(7), and it occurs for a perfectly valid element (4) of the half-Tate
Banach algebra.

This does not invalidate the finite-character criterion of v101.  If the
simple odd characters are critical, their algebraic action extends
abstractly to $\mathcal B_{1/2}$.  What fails is constructing that extension
by applying an arbitrary Banach multiplier to a representative in
$\mathcal H_-$ and removing only the two Tate powers.

Therefore a proof of (8c) must operate on the finite Loewy quotients or on a
new boundary space which contains all the power tails (7) with a positive
contractive norm.  Merely adjoining those tails with their natural Mellin
characters does not solve the problem: the centered generator on
$x^{-s}$ has modulus $n^{1/2-\Re s}$, so positivity plus contractivity again
forces $\Re s\ge1/2$.

\begin{center}
\begin{tabular}{>{\raggedright\arraybackslash}p{.58\linewidth}
                >{\centering\arraybackslash}p{.15\linewidth}
                >{\raggedright\arraybackslash}p{.17\linewidth}}
\toprule
Statement & Status & Location \\
\midrule
Two-oriented Mellin multiplier & \PROVED & (1)--(3) \\
Exact non-Tate residual powers & \PROVED & (4)--(7) \\
Cancellation by the two Tate modes & false & Theorem 1 \\
Abstract Banach extension on critical simple factors & conditional & v101 \\
Contractive boundary realization of all Loewy factors & \OPEN & new theorem required \\
Row (d), condition $G$ & \OPEN & equivalent to the last line \\
\bottomrule
\end{tabular}
\end{center}

\end{document}
